% !TEX root = ../main.tex
\pagestyle{appendix}
\appendix
\section*{Appendices}
\addcontentsline{toc}{section}{Appendices}
\renewcommand{\thesubsection}{\Alph{subsection}}

\subsection{Push Mechanics: Limit Surface and Characteristic Length}
\label{app:mechanics}

The limit-surface normality condition of Equation~\ref{eq:normality} states that
the instantaneous motion twist of a quasi-statically pushed object is normal to
its friction limit surface at the point corresponding to the applied wrench. Its
practical consequence for the task concerns any off-centre push, one whose line of
action does not pass through the object's centre of friction. Such a push produces
a wrench with a non-zero moment component, and therefore a twist that contains an
angular velocity. Almost every push consequently changes both position and
orientation.
The characteristic length $L$ used in the $\mathrm{SE}(2)$ distance of
Equation~\ref{eq:dpose_methods} is the radius of gyration of the object, a
physical constant; for the T-block $L = 0.07\,\text{m}$. It is not a tuning
parameter, and it sets the weight of an orientation error relative to a position
error in a manner consistent with the object's mass distribution. For the
rotationally symmetric disc the characteristic length is set to zero, which
collapses the $\mathrm{SE}(2)$ distance to a pure position distance and encodes the
absence of an orientation objective. The mechanical contrast between the two
objects, the coupled twist of the T-block against the pure translation of the disc,
is discussed in Section~\ref{sec:pushing}.

\begin{figure}[H]
    \centering
    \includegraphics[width=0.6\textwidth]{diagram_limit_surface.png}
    \caption{An off-centre push produces coupled translation and rotation of the
    T-block. The motion twist is normal to the friction limit surface.}
    \label{fig:app_limit_surface}
\end{figure}

\subsection{Training Curves}
\label{app:curves}

Figure~\ref{fig:app_curves} shows the training-time success rate and value loss
for the single-agent and curriculum models. The single-agent model improves
steadily over roughly 3{,}000 iterations, the curriculum model follows a similar
trajectory, and the value loss remains stable throughout, which reflects the
bounded range of the exponential potentials.

\begin{figure}[H]
    \centering
    \includegraphics[width=0.48\textwidth]{cmp_Success_Rate.png}\hfill
    \includegraphics[width=0.48\textwidth]{cmp_Value_Loss.png}
    \caption{Training success rate (left) and value loss (right) for the
    potential-based models.}
    \label{fig:app_curves}
\end{figure}

\subsection{PBRS Parameter Sensitivity}
\label{app:sensitivity}

The potential-based reward uses $k_p = 30$, $k_r = 5$, and $w_{\text{pos}} =
w_{\text{rot}} = 10$, chosen from pilot runs. A systematic grid over these
parameters was not performed within the available budget. The robustness of the
reward to its parameters is therefore not established here. It is recorded as a
limitation in Section~\ref{sec:limitations} and as a direction for future work in
Section~\ref{sec:future_work}. A grid over $k_p \in \{10, 20, 30, 50\}$ and
$w \in \{5, 10, 20\}$, evaluated at the 30-scene gate, would provide the missing
sensitivity analysis.

\subsection{Validation Scene Descriptions}
\label{app:scenes}

The 30-scene held-out suite is defined in \texttt{validation\_configs.py} and
comprises three groups of ten. Scenes 1 to 10 are rotation-dominant, with
orientation changes up to $\pi$ radians and mixed translations. Scenes 11 to 20
are position-only, with the goal orientation equal to the start orientation.
Scenes 21 to 30 combine a position change with an orientation change. Within each
group, difficulty is graded easy, medium, and hard by the start-to-goal distance
and the magnitude of the required rotation. All coordinates are expressed in the
environment-local frame in metres, and orientations in radians.

\subsection{cuRobo IK Pipeline and Workspace}
\label{app:ik}

Each decoded push primitive is executed as a five-phase Cartesian trajectory,
approach, descend, push, retract, and return, spanning 72 physics substeps. Every
waypoint is mapped to arm joint targets by the cuRobo 0.7.5 batched Inverse
Kinematics solver, with the gripper held closed throughout. The reachable
workspace is bounded to the region over the table where the solver returns
elbow-up solutions. Pushes whose start or end points fall outside this region are
clamped to the boundary before execution, so the credit assigned to an action
matches the motion actually performed.

\subsection{Observation Space Specifications}
\label{app:obs}

The single-agent models use a 28-dimensional observation: end-effector pose (6),
object pose and velocity (14), goal pose (6), and the position and orientation
errors (2). The goal-solving agent in the self-play models uses the same
28-dimensional layout, with the goal encoded through the GoalEncoder into an
eight-dimensional latent in the difference variant. The goal-proposing agent uses
a 20-dimensional observation that omits the goal, since it sets the goal rather
than solving it.

\subsection{Confounded Ablation Disaggregation}
\label{app:confound}

The single-agent and self-play models differ on at least eight dimensions at once:
the number of agents (one against two), the goal distribution (fixed random
against adversarial and shifting), the goal encoder (absent against present), the
behavioural-cloning term (absent against present), the historical opponent pool
(absent against present), the rollout length (15 pushes for one agent, against
5 plus 10 split between two), the episode budget (5 pushes for the single agent
against 10 for the goal-solver), and the number of policy updates per iteration
(one
against two). The comparison in this thesis therefore isolates self-play as a
package rather than any single mechanism within it. Four single-dimension
isolations are provided nonetheless: the disc contrast of
Section~\ref{sec:rq2_coupling}, the two Alice reward formulations of
Section~\ref{sec:rq2_reward}, the imitation ablations of
Section~\ref{sec:rq2_abc}, and the goal-encoder ablation of
Section~\ref{sec:rq2_encoder}. The cleanest of these is the reward-formulation
pair, whose members differ only in the goal-proposing agent's reward. It shows no
consistent advantage for either formulation. At 528 environments the
outcome-based variant logs a per-push training rate of 35.8\,\%
($\pm$ 8.8\,pp, 5 seeds) against 32.6\,\% ($\pm$ 0.4\,pp, 5 seeds) for the
time-based variant on the T-block, and both transfer zero held-out scenes. On the
disc the time-based variant leads during training (83.8\,\% against 77.9\,\%
per push,
5 seeds each) yet trails slightly on held-out scenes (6.7\,\% against
7.3\,\%).

\subsection{Training Budget Comparison}
\label{app:budget}

Every Isaac Lab run in the experiment set is budget-matched. The total experience
is held at about 23.8 million pushes, and the iteration count is set inversely
proportional to the number of parallel environments. A sweep therefore changes
the per-update batch size alone. The confound that a larger batch also means
fewer gradient updates is stated wherever a scale result appears. The headline
models
train at 528 environments for approximately 3{,}000 iterations. The
potential-based agent and the improvement-reward baseline are additionally swept
across 256 to 2{,}048 environments, and the two self-play subjects across 256 and
2{,}048. At the 528-environment baseline the potential-based single-agent model
reaches a held-out scene success rate of 79.3\,\% ($\pm$ 3.5\,pp, 5 seeds), and
the curriculum model 80.0\,\% ($\pm$ 4.1\,pp, 5 seeds). The improvement-reward
baseline, trained at up to 2{,}048 environments, is still training at the time
of writing; its held-out validation is \TBD, and until it completes the
environment-efficiency comparison between the two reward families remains open.

\subsection{Secondary Variants}
\label{app:secondary}

The two disc self-play models and the Bob-penalty model are headline subjects,
reported in Table~\ref{tab:main_results} and Sections~\ref{sec:rq1} and
\ref{sec:rq2_coupling}; they are not repeated here. The ablation arms, with the
imitation term removed, with the goal encoder removed, and the disc positive
control, appear in Sections~\ref{sec:rq2_abc} and \ref{sec:rq2_encoder}. Two
groups remain outside the main text. The first is the original, non-simplified
potential formulation at the position-only gate. The second is the
cross-environment set: the single-agent, curriculum, and self-play models
retrained in \texttt{gym-pusht}, one seed each, which solve 3.3\,\%, 9.0\,\%,
and 1.0\,\% of the held-out scenes respectively. Both groups are consistent with
the principal findings: self-play remains far below single-agent training across
every variant.
